\section{Routes that do not decide the location of the zeros, with their exact scope}\label{sec:negative}

This section collects the negative results. Each is a theorem (or, where marked, a numerical observation) about a specific construction as it is defined, and it is stated with the scope in which it is proved.
Three kinds occur.
\begin{enumerate}[label=(\roman*)]
\item \emph{A statement that is false}, with a proof or a counterexample (for example the sharp velocity laws, or freeness detected by zeros).
\item \emph{A construction that is uniform in the zero set}: its conclusions hold for every divisor in the open strip with the stated properties, or for zeta functions of curves over finite fields, or for functions without an Euler product that have off-line zeros. Such a construction cannot by itself distinguish RH from its failure.
\item \emph{A construction equivalent to RH by construction}: it isolates the off-line part of the zero set in its definition, so that it vanishes (or is injective, or positive) exactly when there are no off-line zeros. Such a statement is an exact reformulation; it does not supply a mechanism that forces the off-line part to vanish.
\end{enumerate}
None of these results says that the construction concerned is irrelevant to the Riemann hypothesis. Several of them identify an input that the construction lacks: for the weight routes, a pole bound uniform in the tensor degree or a discreteness input with a square (on the programme's finite data the first is equivalent to $\max\re\rho\le\tfrac12$, so it would have to come from an independent source); for the mirror and two-line routes, a property that distinguishes $\zeta$ from the Davenport--Heilbronn function, such as the Euler product. Item N-Prim identifies what the primary decomposition needs: polynomial bounds on the principal parts of $1/\zeta$ at the zeros. Whether the programme's constructions can supply these inputs is open. The numbers in brackets are those of the register's negative results (\note{00\_RESULTS\_REGISTER.md}, Goal~1).

\subsection{Statements that fail}

\negres{N-Vel (sharp velocity laws) [1].}{Under (H), all but finitely many zeros simple and on the line, $\sum_{\gamma\le T}\re c(\rho)=T/4\pi+O(T)$ is false: $\limsup|\re E(T)|/T=\infty$; the vertical law's error is not $O(T)$ either. Unconditionally only the smoothed laws are proved, and for fixed smoothing width their error is of order $T$. Whether the observed skew ($\re E/T$ between $-3.15$ and $+0.17$ up to $T=64\,000$) persists is open, and the one-sided statements $\limsup\re E/T=+\infty$ and $\liminf\re E/T=-\infty$ are not proved. (Copy, Theorem~C; proof completed in \note{20\_}.)}

\negres{N-Sheet (the non-Eulerian sheets) [3].}{For $a\in(0,1)\setminus\{\tfrac12\}$, $\zeta(s,a)$ has infinitely many zeros with $\re s>1$ (Davenport--Heilbronn \cite{DavenportHeilbronn} for rational or transcendental $a$; Cassels \cite{Cassels} for algebraic irrational $a$). For $a>1$ such zeros are found numerically; at $a=2$, $\zeta(s,2)=\zeta(s)-1$ vanishes at $1.40778804065+23.3279877546i$. So an argument that uses only properties shared by all sheets $\zeta(s,a)$, $0<a<1$, cannot prove that the nontrivial zeros of $\zeta$ lie on the line; it must use a property of the Eulerian sheets $a\in\{\tfrac12,1\}$ (Corollary~\ref{cor:eulerian}), such as the Euler product.}

\negres{N-Beurling (regularity of the integers) [12, 13].}{Regularity alone does not force RH for generalised primes: a Beurling system can have $N(x)=kx+O(x^{1/2}\exp(c(\log x)^{2/3}))$ and zeros approaching $\sigma=1$ (Zhang \cite{Zhang}; also Diamond--Montgomery--Vorhauer \cite{DMV}). Near-exact counts do not pin the prime system either: $(\text{primes}\setminus\{2\})\cup\{3',4\}$ (a second copy of $3$) has $N(x)=x+O(\log^2x)$ (copy, T2). Exact counts do pin it (the first-appearance sieve). So the question the programme poses is how the additive counting step constrains the multiplicative system, since regularity alone does not.}

\negres{N-Free (zeros do not detect freeness in general) [28].}{The monoid $\{1\}\cup\{2^e:e\ge2\}$ is not free ($64=4^3=8^2$), yet its zeta function $(1-x+x^2)/(1-x)$, $x=2^{-s}$, has all its zeros on $\re s=0$. For congruence monoids ``no zeros in $\re s>1$'' does imply freeness (Proposition~\ref{prop:congruence}, by Saias--Weingartner); the formal-logarithm test (Theorem~\ref{thm:formallog}) works for every monoid of integers.}

\negres{N-Eu (finitely many Euler factors) [50].}{No finite set of primes fixed in advance has Euler factors that determine every finite shifted-zeta product. For primes $p_1,\dots,p_k$ and $a_j=2\pi i/\log p_j$, the even and the odd sub-sums of the box $\{\sum_j\varepsilon_ja_j:\varepsilon\in\{0,1\}^k\}$ give products of shifted zeta functions with identical Euler factors at every $p_j$ but different exponent measures (their difference has Laplace transform $\prod_j(1-e^{a_jz})\ne0$). So Lemma~\ref{lem:shifted} and its corollary are global statements; for a particular product finitely many factors can suffice. (claude-ab, \note{29\_} 29.N2.)}

\negres{N-Mirror (the mirror line at $-\tfrac12$) [37].}{The author's second line is exactly $A(\Zs)=\Zs-1$, unconditionally, and ``mirror $=$ shift'' pointwise is equivalent to RH (Theorem~\ref{thm:mirror}). The properties used there (symmetry of the zeros under $s\mapsto1-\bar s$, their location in the strip, and an explicit formula with a positivity criterion) all have counterparts for zeta functions of curves over finite fields, where RH is Weil's theorem \cite{Weil1948,Rosen}, so they cannot imply an off-line zero. The $0$-centred partner pairing is not Hermitian (Proposition~\ref{prop:centres}(b)). Coexistence of the two lines does not decide either: the statements for $\zeta(s)$ and for $\zeta(s+1)$ are one proposition, and the Davenport--Heilbronn function coexists with its translate, has a functional equation of $\zeta$'s kind, and has the zero $0.808517182456637+85.6993484853776i$ off the line; what it lacks is an Euler product. (\note{17\_} §§5--6, 10.)}

\negres{N-Logic (independence) [43].}{RH cannot be independent of ZFC and false, since it is equivalent to a $\Pi_1$ sentence (Proposition~\ref{prop:detect}). This gives no shortcut: non-refutability is equivalent to RH for $\Sigma_1$-sound theories.}

\subsection{Positivity and receivers see only the critical zeros}

The programme builds many positive forms, Hilbert completions and receivers on the zeta quotient. The exact results form one picture: a positive construction compatible with the scaling transfer sees only critical zeros (a theorem: the classifications of Section~\ref{ssec:receivers}); each of the constructions listed below that sees the off-line zeros is zero, or injective, or positive, exactly under RH; and the residue duality $B_\zeta$, which is nondegenerate unconditionally, is not positive.

\negres{N-Pos1 (positive transfer forms) [6, 26].}{Every continuous positive form compatible with the scaling transfer vanishes on the off-line part $R=Q/N_O$ (\lbl{CPS4.2}); the classification of such forms (register S12, S22, S35, S48; Lemma~\ref{lem:PTQ}) makes this exact. \lbl{CFP}, \lbl{CPS}, \lbl{TWC10} and \lbl{SPF9} are related results, not independent ones; for tensor degree $k\ge2$ reflected pairs survive, with modulus exactly $n$. So an argument through such receivers cannot remove $R$ by itself.}

\negres{N-Pos2 (the $L^2$ norm, harmonic sweeping, Hilbert quotients) [24, 25, 30, 33, 35].}{The summation image is dense in $L^2(du)$ by Wiener's $L^2$ Tauberian theorem, so the naive $L^2$ quotient seminorm is zero (\lbl{ASD10}, \lbl{OPD3}). Poisson-sweeping the zeros onto the line gives a positive distribution that descends to $Q$ if and only if there are no off-line zeros (\lbl{HSW6}). The Hilbert receivers are automatically pure (prime spectra on $|z|=\sqrt p$, unconditionally), and the kernel of the map to their inverse limit is exactly the off-line jet classes, zero if and only if RH holds (\lbl{WHR8--WHR12}); a spectral-radius argument there proves purity of what remains after the loss, and cannot by itself prove RH (an assessment). A seminorm invariant under $T_p$ and $T_p^{-1}$ vanishes on every primary block with $0<\re\rho<1$, and with the half-density action it vanishes on every off-line block (\lbl{ATG6.5}; \note{14\_} §2). The off-line functional $\Theta_{\mathrm{off}}=\sum_{\re\rho\ne1/2}m_\rho F(\rho)$, continuous on $Q$ and zero exactly when there are no off-line zeros, is never a nonzero multiple of $b\mapsto b(1)$: an explicit $b_*$ in the summation image has $b_*(1)=2$ and $\Theta_{\mathrm{off}}(b_*)=0$ (\lbl{VWR10.5--10.9}).}

\negres{N-Pos3 (RH-equivalent objects) [47, 48, 53, 54, 56].}{The following are equivalent to RH, or to RH with simplicity, by construction, each saying only that there are no off-line zeros (or no multiple zeros): \lbl{AST}'s $\mathcal{R}$, \lbl{GDC}'s $N_O$, \lbl{PTQ}'s $H_{\mathrm{off}}$ and the descent of the Weil form, the traces and detectors of \lbl{PTQ8.3} and \lbl{ECR5.5} (eight constructions, three conditions); $W\ge0$, $W$ positive definite (with simplicity), $U_a=D_a^*$ for one $a\ne1$, density of the trace image (simplicity) (\lbl{RTT}, \lbl{RZ}); the ordinary trace of \lbl{GDE}'s energy, normality of $B_{r,t}$, the vanishing of the degree-adjusted trace, of \lbl{ABH}'s $\mathcal{L}_{r,t}$, of \lbl{HSR}'s $D_{r,t}$ and $\Xi$, of \lbl{HBW}'s singularity-class map, unitarity of $n^{-1/2}A_n$, weight-$\tfrac12$ counting moduli, monodromy phases $-1$, and $D_{r,t}\subset H^2$. The two numerical detectors are not holomorphic in $\rho$, so the explicit formula does not compute them directly. Every attempted cancellation (cyclicity, the degree relation, the cover relation) leaves the off-line sector intact, and the unconditional results among them (\lbl{NHI5.4}, \lbl{NHI8.4}, \lbl{HSR4.1}, \lbl{HSR5}, \lbl{HSR6}, \lbl{HSR10.6}, \lbl{HBW3}, \lbl{HBW8}, \lbl{HBW9}) hold whether or not RH does. The adjoint-domain route on the Hardy boundary receiver is circular: $\mathcal{L}_{r,t}x\in\mathcal{D}_{\delta_1}$ if and only if $x=0$ (Noor's unconditional theorem). (\note{25\_} §5, \note{27\_} §2, \note{30\_}, \note{32\_}.) The zero-set condition $D_{r,t_2}=D_{r,t_1}$ (finitely many off-line zeros) is weaker than RH, and the two conditions read off the zero-time Gaussian energy (compactness, and the trace-class condition $\sum m_\rho(\beta-\tfrac12)^2<\infty$) are implied by RH (register S58); their status is not claimed.}

\negres{N-Pos4 (the divisor potential) [29].}{Under $s\mapsto1-\bar s$ the odd part of $\log|\zeta|$ is exactly $\tfrac12\log|\chi|$, and its contribution to the off-line detector $\mathcal{E}_r(t)$ is accounted for by the pole and the trivial zeros; the detector lives entirely in the even part, about which the functional equation says nothing further (\lbl{OZD5}, \lbl{OZD7}). For fixed $(r,t)$ the Gaussian weight $e^{-2t\gamma^2}$ makes each such identity numerically blind above modest heights (an assessment).}

\subsection{Constructions that are uniform in the zero set}

\negres{N-Unif1 (completeness, integrality, clocks) [14, 15, 16].}{``Off-critical, therefore undetected'' is false: the test function $\varphi(v)e^{-(\rho-1/2)v}$ detects any zero, and an off-critical quartet would appear in the absorption measure as $-2m(e^{\beta t}+e^{(1-\beta)t})\cos\gamma t$, which completeness supplies no sign or growth bound to exclude (\lbl{IH2--IH4}). Clock changes preserving order and composition are linear and move the critical line with the zeros (\lbl{PL2--PL8}). No finite-rank integral model of the dilations exists (\lbl{IH5.3}).}

\negres{N-Unif2 (lifting and the restriction row) [18, 19, 20, 22].}{No eigenvector maps onto the source character at any zero; the least lift of the $k$-th jet has order $m+k+1$, which depends on the multiplicity, not on the position (\lbl{MCL6}). In the $\zeta$ restriction row, kernel, middle term and target carry the same character $1-\rho$, so any weight truncation keeps or kills the whole extension, and the fixed-order class is nonzero at every zero, on or off the line (\lbl{ORE5.8}, \lbl{ORE6.5}); the residue pushout removes the obstruction at every zero (\lbl{RPC6.3}). \lbl{MCL3--MCL7}, \lbl{ORE6.2}, \lbl{ORE6.5}, \lbl{IH6.4}, the first two laws of \lbl{IH6.5} and \lbl{PL2.1} hold word for word for the even lattice sum at $a=\tfrac15$, whose factorization is not unique and which has zeros in $\re s>\tfrac12$ (\note{08\_} §2).}

\negres{N-Unif3 (covers and sections) [36, 55].}{For the transposed cover action on the dual of $\B+\C h_{P,t}$, the largest cover-invariant part of the zero-coordinate slice is exactly the divisor that the chosen residue section inserts: the zeta divisor for $\xi$, any finite prescribed set of nonzero points for a polynomial section, nothing for the constant section (\lbl{RSS3.2--RSS3.3}, \lbl{RSS6}). Two residue sections of the same non-split extension have cocycle spans $I_\zeta$ and $\B$ (\lbl{NJS10}). Covariance gives recovery, not location.}

\negres{N-Unif4 (the two-line splitting and the second copy) [38, 44, 46].}{The splitting $\B/(I\cap I_-)\cong\B/I\oplus\B/I_-$ and the derived separation of $Q$ and $Q_+$ hold for every primitive Dirichlet $L$-function (Proposition~\ref{prop:separatorL}); their proofs never use where the zeros lie inside the open strip, and the Euler product enters only through zero-freeness of $\re s\ge1$. The derived separation is the general vanishing of Ext for coprime annihilators. $\B$ itself does not split along the two lines. The statements about the copies that involve RH restate it through $\mathcal{R}=Q/N_O$, which records only values at off-line zeros; the programme itself says that no vanishing and no new pure weight follows (\lbl{ADM9}).}

\negres{N-Unif5 (returns, windings and monodromy) [58, 59].}{\lbl{BML}'s winding difference is invertible for every zero set avoiding the integers with $e^{-2\pi i\rho}$ at positive distance from $1$ (Lemma~\ref{lem:winding}), and \lbl{BML11}'s $H^1$ is an artifact of truncating the logarithmic tower. The iterated-return and reflected-return theorems of \lbl{BRC}, \lbl{GJN} and \lbl{GJNR} hold for every divisor in the open strip with their other hypotheses (Lemma~\ref{lem:escape}); the norm-shift relation $NP_b=bP_bN$ is the chain rule and has solutions for every nilpotent.}

\negres{N-Unif6 (the doubled source, the adelic complex, the jet map) [69, 70, 71, 72].}{The programme's doubled source adds nothing to Connes--Consani's cohomology (Section~\ref{sec:other}); the source-space gluing and the doubled pullback (\lbl{CGS4--CGS10}, \lbl{DCP4--DCP12}) are uniform in the zero set, which enters only as the spectrum of the dilations on $H^1$ and through the jets: in finite models with zeros on and off the line all dimensions, ranks, signs and sections agree. The lift in \lbl{DCP} is Connes--Consani's step, which vanishes by Poisson summation, not by a separation of weights as in Weil~II~§3.6 (\lbl{DCP12} says the comparison has not been assumed). The Ext-vanishing between the zeta quotient and the endpoint modules sees only $\xi(0)=\xi(1)=\tfrac12\ne0$ (\lbl{GEX9} states the limitation). \lbl{OMS5--OMS7A} use only $0<\re\rho<1$, the multiplicities and the unboundedness of the heights (\lbl{OMS8} says so).}

\negres{N-Prim (the primary decomposition) [73].}{The primary decomposition of the zeta quotient is not available without new input about $\zeta$ at its zeros: its convergence for every class is equivalent to polynomial bounds on the principal parts of $1/\zeta$ at the zeros, and for simple zeros to $|\zeta'(\rho)|\ge c(1+|\gamma|)^{-K}$, which would imply that all zeros are simple (Theorem~\ref{thm:primary}). $Q$ is not the product of its primary blocks (Corollary~\ref{cor:jetimage}); the constant tuple $1$, the obstruction behind the programme's nonzero class $c_\zeta$ (\lbl{PGD4}), is not a jet tuple.}

\subsection{Weights, tensor powers and the square}

\negres{N-W1 (the weight separations proved so far) [17, 21, 65].}{The programme's receivers separate weights using only the strip bounds $0<\re\rho<1$, or Tate shifts built into the receiver; a geometric source for a dividing line at weight~$1$ has not been constructed (\lbl{TL8}; an assessment). Positivity of the timed-prime measure reaches exactly $\re s=1$: it gives $\zeta(1+it)\ne0$ and, with standard growth bounds, the classical zero-free region, and no zero-free strip (Beurling systems with positive prime measure can have zeros approaching $\re s=1$). The method of Weil~II~§2.1, specialised to $\zeta$ (Deligne's example 2.1.9), proves only these classical statements; the compact form, the cohomological proof of condition (C), the strict bound and Weil~II~§3 have no $\zeta$ analogue along this route (\note{37\_} 37.N2--37.N3).}

\negres{N-W2 (tensor powers and the square, as constructed) [23, 27, 31, 32, 34, 52, 66].}{Positivity of even tensor powers returns $\max\re\rho$ and cannot bound it (Proposition~\ref{prop:abscissa}); the missing input is a pole bound at $k+1+C$ uniform in $k$, which on the programme's finite data is equivalent to $\max\re\rho\le\tfrac12$. The characters of source and target match in every tensor power; the transfer-compatible positive forms kill every tuple off the centred set; the tensor Weil form is indefinite on the reflected sector; the square's obstruction vanishes exactly when the original one does (\lbl{TWC6}, \lbl{TWC9--TWC11}, \lbl{FST3}, \lbl{FST4.2}, \lbl{FST4.4}, \lbl{FST7}). In degree~$1$ the square adds nothing beyond $Q\otimes Q$ (\lbl{GMC4--GMC6}), and reflected off-line pairs are centred ($\rho+\rho^\#=1+2i\gamma$). On the Hilbert-closure defect, source and target of the connecting map have the same character (\lbl{HCS7}, \lbl{HCS10--HCS11}). The faithful norm on $Q^{\otimes k}$ is bounded with exponent $k$ (sharpness is proved only for the source seminorms, \lbl{ATG6.1}), and the critical Hilbert components have exponent $k/2$ but common kernel exactly the tensors with an off-line factor; for finite reflection-stable $S$ the exact finite-trace exponent is $k\max\re\rho$, so the error against $k/2$ is linear in $k$ (\lbl{ATG}). The tensor Euler products of an off-line quartet have their genuine pole at $1+2k\max(\re\rho,1-\re\rho)$ for every complex $\rho$, zero of $\zeta$ or not (checked at $\rho=0.7+10i$), so they cannot detect zeros (\note{37\_} 37.N5).}

\negres{N-W3 (the average weight and determinant data) [60].}{The average weight of an off-line quartet is exactly $1$ (Proposition~\ref{prop:avgweight}); only the constituent weights see an off-line zero, and Weil~II~1.5.1 turns a pure determinant weight into purity but cannot produce one. Determinant data at a point control only the sum of the weights (\note{36\_} 36.N1).}

\negres{N-Tr (the product trace and the diagonal return) [39].}{The product-trace machinery does not force weight~$1$. The reflected-pair classes have character $p$ (weight~$2$) whether or not either zero is on the critical line, and the full global complex kills them with an explicit boundary (\lbl{FEM7--FEM8}). The diagonal return keeps the product trace exactly but pairs the weights $2\re\rho$ and $2-2\re\rho$ (\lbl{DER10}); in the programme's words it does not prove either individual weight to be one. Also exact: no degree-zero same-site lift exists (\lbl{FTD9}), no continuous single-valued lift of $q$ fixes the arithmetic diagonal (\lbl{FSC2}), and the comparison $R\Delta^!K\to\Delta^{-1}K$ is zero (\lbl{DER6}) (\note{19\_}).}

\negres{N-W4 (Deligne's local theory against the programme's operators) [61, 62, 63, 51].}{The boundary bounds on $V$ and $V^\vee$ alone do not force the conclusion of Weil~II~1.8.4; the tensor square is needed (Proposition~\ref{prop:purity}). An operator commuting with $N$ cannot carry Deligne's local weights unless $N=0$, which applies to the programme's counting operators $W_p$ and $C_b$ on jet blocks of length $\ge2$; the ramification operator $P_b$ has no eigenvector on its space, so the weight filtration that Deligne builds from a Frobenius is not defined for it, and only the relation $NP_b=bP_bN$ transfers. Tensor Euler products are blind to jets, and the programme's residue extension is exactly such a non-split Jordan block, at the pole $s=0$ (\lbl{ETR10}, \lbl{NPE4}).}

\negres{N-W5 (the recovered arithmetic base) [67, 68].}{The programme's recovered ring $R=\operatorname{End}(G/G_{\mathrm{tors}})\cong\ZZ$ carries only the free rank of the winding group; every automorphism acts trivially on it, so $\operatorname{Spec}R$ is $\operatorname{Spec}\ZZ$ with its standard arithmetic, and the torsion, the chart twist, the involution and the support $\tau$ do not reach it. Every Frobenius eigenvalue of the objects built there from programme data (the constant sheaf, its Tate twists, $i_{p*}E$, the costalks at closed points of $\operatorname{Spec}\ZZ[1/\ell]$) is an integral power of the residue characteristic, so all these weights are even; the purity circle $|a|^2=p$ needs weight~$1$, which Deligne's theory supplies on this base through, for example, $H^1$ of an elliptic curve over $\mathbb{F}_p$, but which no object built from programme data supplies (\note{38\_}).}

\negres{N-W6 (finer points of Weil~II~§2) [64].}{Positivity alone cannot exclude Deligne's quadratic exception (Lemma~\ref{lem:intpos}); in Weil~II the double cover does, and for $\zeta$ the absence of a nontrivial quadratic character of $\R$.}

\negres{N-F1 (items of the $\mathbb{F}_1$ lane) [7, 8, 9, 10, 11].}{Doubling has no sign-preserving mirror-compatible lift on Connes--Consani's signed data (a lift killing the sign does exist). The compact-completion modulus theorem holds for every compact group (\lbl{P12} says so). The $\sqrt n$ dilation radius cancels in the product formula. The twistor involution does not fix the radius (equal radii only on $|z|=1$). The weight-one sign sector is $L(E,s)$ of conductor $32$, not $\zeta$'s.}

\subsection{The two-line system}

\negres{N-Two1 (parity and the four-direction anomaly) [40].}{The $\ZZ_2$ anomaly of one line counts only on-line zeros; in the two-line system the $\ZZ_2$ anomaly vanishes identically and the $\ZZ_4$ anomaly equals the one-line parity; the integer count sees off-line pairs only together with $N_0(T)$, and the holonomy index does not separate off-line pairs from multiple zeros on the line (Propositions~\ref{prop:parity} and~\ref{prop:index}).}

\negres{N-Two2 (the structure without the Euler product) [41, 45].}{The Davenport--Heilbronn function has the completed functional equation, the reality, the nonnegativity of its two-line product on $\re s=0$ (strict positivity is not established for it), and a size-four orbit for its off-line zero; the superposition of $L(s,\chi)$ and $L(s,\bar\chi)$ cancels the root-number phase by addition and loses the Euler product, while the product $L(s,\chi)L(s,\bar\chi)$ cancels it by multiplication and keeps it (\note{22\_} 22.10--22.11). The separator proof does not reach the Davenport--Heilbronn function (Section~\ref{sec:twolines}); whether the splitting fails for it is open.}

\negres{N-Two3 (the second line on its own) [74].}{For products $\Lambda$ of completed $L$-functions of primitive Dirichlet characters in which $\chi$ and $\bar\chi$ occur equally often, the positivity $G_\Lambda(it)=|\Lambda(1+it)|^2>0$ follows from $\Lambda$'s functional equation, reality and nonvanishing on $\re s=1$, and $G_\Lambda$ and its explicit formula are functions of $\Lambda$, so the combination adds no input to the count of off-line zeros (this does not decide whether one-line inputs can bound that count better than Turing's method). The chiral cross form is hyperbolic whether or not RH holds; only its diagonal part, which is Weil's form, sees off-line zeros (Proposition~\ref{prop:chiral}). The two-line prime race differs from the one-line race by a convergent series, so its limiting distribution (under GRH) and its logarithmic density (under GRH and linear independence) are unchanged (Proposition~\ref{prop:races}).}

\subsection{Cross-programme transfers}

\negres{N-X (transfers between the workbenches) [42, 49, 57].}{Joint averages do not survive restriction of the Navier--Stokes torus cover to unchanged finite torsion ($13$ against $4$ at $p=13$, $a=4$); the complete preimage repairs this. The Navier--Stokes profile gives Yang--Mills a zero-quotient sequence only as the coupling tends to $0$, and only as one admissible weight. The Jacobian $\to$ Yang--Mills chain yields no fixed-coupling or continuum statement: the Yang--Mills workbench's Theorem~14.2 is correct conditional on standard fixed-box weak-coupling limits that the file only sketches, and its energy quotients are, up to $1\pm1/(10j)$, the lowest free two-gluon colour-singlet energy of an open box of side $j/50$ at couplings $g_j<1/j$, so $Q_j\cdot(\text{box side})\to2\sqrt2\pi$, the scaling of a massless free field, whether or not the continuum theory has a gap; the Jacobian tensor enters the norm and the energy, not the limit of the quotient (\note{34\_}). The $S^6\to$ Yang--Mills edge uses only $D>0$ on its patch; the Navier--Stokes workbench certifies none of its early heat or $S^6$ transfers; prime-only clock seeds are not transfer-stable (\note{21\_}, \note{28\_}).}
